% RESEARCH_DRAFT_SHA256: 8d3a397e838bbb65681b623fb8edee314b03ee0dce3e993209f7f892ad57926f
% Layout and source-reference wrappers expanded for the web renderer.
\documentclass[11pt]{article}
\usepackage[margin=25mm]{geometry}
\usepackage{amsthm}
\usepackage{fontspec,unicode-math}
\setmainfont{Latin Modern Roman}
\setmathfont{Latin Modern Math}
\usepackage{amsmath,amssymb,mathtools}
\usepackage{xurl,hyperref}
\hypersetup{colorlinks=true,urlcolor=blue,linkcolor=blue}
\setcounter{secnumdepth}{0}
\setlength{\parindent}{0pt}
\setlength{\parskip}{5pt}
\setlength{\emergencystretch}{3em}
\newcommand{\sref}[2]{\hyperlink{source:#1:#2}{[S#2]}}
\newcommand{\cataloguescope}[1]{\paragraph{Source question.}#1}
\newcommand{\floor}[1]{\left\lfloor#1\right\rfloor}
\newtheorem{proposition}{Proposition}
\newtheorem{lemma}[proposition]{Lemma}
\begin{document}
\paragraph{Research manuscript.} OpenAI \texttt{gpt-6-astra}, reasoning \texttt{ultra}; 11 October 2026.
\paragraph{Preserved source body.} The complete catalogue mathematical text and sources follow with layout and source links expanded, including source qualifications and its historical preparation notes. New research begins at the section headed Research attempt.

% UnsolvedMath ID 30005828; source catalogue code OWR-14298166-004
\section{Superpolynomial Lower Bounds for Dag-Like Resolution over Parities}\label{30005828}
{\small UnsolvedMath ID 30005828 · Computer Science}
\subsection{Definitions and mathematical statement}

\paragraph{Shared notation.}Standard number systems are denoted by $\mathbb N,\mathbb Z,\mathbb Q,\mathbb R,\mathbb C$; index ranges, computational models and specialized notation follow the cited problem source and the definitions below.


An affine parity equation is $\sum_{i\in S}x_i=b$ over $\mathbb F_2$, with $b\in\{0,1\}$. A parity clause is a disjunction of such equations. Resolution over parities, $\operatorname{Res}(\oplus)$, permits resolving complementary parity equations and the source linear-algebra inference rules. A refutation derives the empty clause from an unsatisfiable formula. Dag-like proofs may reuse previously derived clauses, whereas tree-like proofs duplicate them when reused. Size measures the encoded derivation under the source convention. A superpolynomial lower bound exceeds every fixed power of the input size for a suitable growing family of unsatisfiable instances.
\paragraph{Statement recorded in the supplied source.}
Prove a superpolynomial lower bound on the size of dag-like $\operatorname{Res}(\oplus)$ refutations. \href{https://doi.org/10.4171/owr/2024/15}{[S2]}
\paragraph{Scope qualification.}
The target is the general dag-like system. Lower bounds for regular, tree-like, or otherwise restricted parity resolution do not settle it.
\paragraph{Source status.}
The supplied catalogue labels this record open. Its source-status metadata was checked against the supplied snapshot on 2026-10-11. This import does not independently establish current open status or mathematical correctness. The catalogue literature assessment is dated 2026-08-22.
\subsection{Short English statement}
Find formulas for which every unrestricted parity-resolution refutation is larger than any polynomial in the input length.
\subsection{Sources}
\begin{itemize}
\item[S1] Ulam AI and the UnsolvedMath Contributors, UnsolvedMath: A Curated Collection of Open Mathematics Problems, record 30005828 (OWR-14298166-004). Catalogue curation and metadata: CC BY 4.0. The supplied source question is retained; definitions and the English summary are editorial additions.
\url{https://huggingface.co/datasets/ulamai/UnsolvedMath}.
\item[S2] Proof Complexity and Beyond (2024), Oberwolfach Reports, DOI 10.4171/owr/2024/15, pp. 871--934
\url{https://doi.org/10.4171/owr/2024/15}.
\item[S3] Efremenko, Garlik, and Itsykson, Lower Bounds for Regular Resolution Over Parities
\url{https://doi.org/10.1145/3727978}.
\item[S4] Oberwolfach Report 15/2024
\url{https://ems.press/content/serial-article-files/48654}.
\textit{Underlying Oberwolfach Reports material retains its source-specific licence, including CC BY-SA 4.0 where applicable; see THIRD\_PARTY\_NOTICES.md.}
\end{itemize}
\paragraph{Preparation and review.}
Original question attributed to the cited source and UnsolvedMath contributors. Codex (GPT-6; exact serving revision not exposed) assisted with definitions, English summaries, source inspection and formatting. No new human mathematical authorship or independent proof verification is asserted. Canonical ID 30005828 and source code OWR-14298166-004 checked against the complete supplied registry and existing repository statements; known flagged repetitions and obvious identical conjectures excluded. Full mathematical-equivalence/literature review has not been performed. Statement preparation and final review followed contribution prompts C and D; automated checks verify structure and identity, not open status.
\label{end:30005828}

\section{Research attempt}
\subsection{Literature, source context and scope}
The full original catalogue body and source qualifications are preserved above. The relevant primary workshop passages were read in the cached publisher report. Historical claims and subsequent attributed results are distinguished from the mathematical work below. A restricted example, standard lemma or known theorem is not counted as new progress on the entire original target. Missing assumptions and repeated or rhetorical extracts are not silently turned into different solved problems.

\subsection{Testing the pure parity-cycle route}
The primary source distinguishes regular resolution over parities from the general system. The 2026 Itsykson--Knop paper likewise states the general size question as open while establishing size--depth tradeoffs. We use its semantic weakening convention: from a linear clause one may infer any linear clause it implies, and resolve complementary equations $f=0,f=1$. A depth tradeoff alone is not an unrestricted size lower bound.
Consider $m\ge3$ variables on a cycle and unit parity clauses
\[
x_i+x_{i+1}=b_i\quad(1\le i\le m),\qquad x_{m+1}=x_1,
\qquad \sum_i b_i=1
\]
over $\mathbb F_2$. These are already parity clauses, not an unspecified CNF encoding.
\begin{proposition}
This contradictory family has a tree-like resolution-over-parities refutation with $O(m)$ lines and $O(m^2)$ bits in an uncompressed dense representation.
\end{proposition}
\begin{proof}
From units $f=b$ and $g=c$, weaken the second to
\[
(f=1-b)\ \vee\ (f+g=b+c).
\]
Every assignment satisfying $g=c$ satisfies this disjunction: if $f\ne b$ the first disjunct holds, and otherwise the second does. Resolve with $f=b$ to infer the unit $f+g=b+c$. Repeatedly applying this operation to the cycle equations yields after $j<m$ equations the unit $x_1+x_{j+1}=\sum_{i\le j}b_i$. The last addition gives $0=1$, which semantically implies the empty clause. Each earlier derived unit is used once, each input is used once, and each addition costs two lines. This proves the tree and line bounds. A linear equation on $m$ variables uses at most $O(m)$ bits, so even dense encodings have the stated total size.
\end{proof}
Thus an attempted lower bound based merely on global inconsistency of local parity constraints fails: this system can accumulate that inconsistency directly. This does not say all Tseitin encodings are easy in all proof systems, or identify a hard family for the general parity system.
\subsection{Verification and remaining gaps}
All assignments for cycle sizes 3 through 10 were checked for unsatisfiability; the addition inference was truth-table checked over its possible parity values. The derived prefixes and line counts were also verified. These checks concern the specified unit-parity inputs. No lower bound for arbitrary DAG-like derivations follows, and a regularity or depth restriction cannot be omitted from an external theorem.

\section{References and source attribution}
Primary workshop report: \url{https://doi.org/10.4171/owr/2024/15}. Relevant passages read: 2024/15 p.897, regular resolution over parities and the general-system question.

\url{https://drops.dagstuhl.de/storage/00lipics/lipics-vol362-itcs2026/html/LIPIcs.ITCS.2026.81/LIPIcs.ITCS.2026.81.html}. Read Introduction inference rules, open general-size statement and Theorems1--2 describing size-depth tradeoffs; no full external proof verification.

All retrievals were checked on 11 October 2026. Original question attribution remains in the preserved Sources section. External theorems are attributed and not claimed as new discoveries or fully verified merely by citation.


\paragraph{Integration review.} Codex (GPT-6; exact serving revision not exposed) reviewed the written arguments and expanded layout and source-reference wrappers. Repository metadata and the production web renderer were checked using Python 3.13 and JavaScript with MathJax 3.2.2. This model-assisted review is not independent human proof verification; no PDF compilation or formal proof checking is asserted.

\section{Completion Estimate}
\noindent\textbf{COMPLETION: 0\%}
No substantiated new progress toward the entire original target is claimed. The worked results and failed route are delimited above; a known result, illustrative model, or auxiliary obstruction is not counted as original-target completion.

\section{Final status and attribution}
\textbf{Attempt status: unresolved.} This describes this research attempt, not a new ruling on the source catalogue's status or on external literature claims.
The mathematics and this self-audit were generated by OpenAI \texttt{gpt-6-astra}, reasoning \texttt{ultra}, on 11 October 2026. Codex (GPT-6; exact serving revision not exposed) orchestrated the batch, read and reviewed its written mathematical arguments, and contributed formatting/repository review. No human mathematical authorship, independent human proof checking, or formal verification is claimed. The historical preparation paragraphs in the reproduced catalogue body concern that body only.

Contribution prompt A was applied to the complete definition and supplied source record, followed by the adversarial checks of prompt B. Local reproduction files are \texttt{batch025/checks.py}, \texttt{computations.json}, and the per-ID audit JSON. They were not public links at preparation time. Numerical checks are finite and do not prove unrestricted results. No LaTeX engine was found locally; compilation was not performed by this research agent.

\end{document}
